% Part: proof-theory
% Chapter: proof-search
% Section: tableaux

\documentclass[../../../include/open-logic-section]{subfiles}

\begin{document}

\olfileid[tr]{pt}{ps}{tab}

\olsection{Tableau Sistemleri}

Tableau sistemleri, sekant hesaplarıyla yakından ilişkili olan ve hem
elle hem de yazılımda uygulanarak ispat aramaya özellikle elverişli
!!{proof} sistemleridir. Tableau sistemlerinde sekantlardan oluşan ağaçlar
kurmak yerine, !!a{proof}, !!{formula} biçimlerinden oluşan bir ağaçtır. Ancak
doğal tümdengelimden farklı olarak kurallar sekant hesabının kurallarını
yansıtır. Bir tableau !!{proof}, özünde açık bir model aramasıdır. Bu
nedenle bu sistemlere bazen \emph{anlamsal} tableau ya da \emph{doğruluk
ağacı} da denir.

\emph{İşaretli} tableau, $\sFmla{\True}{!A}$ veya
$\sFmla{\False}{!A}$ biçiminde olabilen işaretli !!{formula} biçimlerinden
oluşan bir ağaçtır. Ağaçlar aşağı doğru büyür. İstenen sonucu belirten
!!{formula} biçimleriyle başlarlar; amaç bu sonucu !!{prove}{}maktır. Örneğin,
$!A, !B \Sequent !C, !D$ sekantını !!{prove}{}mak için tableau
$\sFmla{\True}{!A}, \sFmla{\True}{!B},
\sFmla{\False}{!C}, \sFmla{\False}{!D}$ ile başlar.
$!A$ ve $!B$'yi doğru, $!C$ ve $!D$'yi yanlış yapan !!^a{structure},
$!A, !B \Sequent !C, !D$ sekantının geçerli olmadığını gösteren
!!a{structure} olur.

Tableau ağacının bir yaprak (en alttaki) düğümü, aynı dala ek düğümler
ekleyen ya da aşağıya iki kardeş düğüm ekleyerek dalı iki yeni dala
ayıran dal genişletme kurallarıyla genişletilebilir. Kurallar, aynı
dalda yukarıda bulunan herhangi bir işaretli !!{formula} için
uygulanabilir. Klasik tableau sistemi \Log{Tc}'nin dal genişletme
kuralları \olref{tab:Tc}'de verilmiştir.

\subfile{rules-Tc}

Görüldüğü gibi, kurallar yalnızca \Log{G3c}'nin baş aşağı çevrilmiş
kurallarıdır; $\True$ ve $\False$ işaretleri ise sırasıyla bir
ana !!a{formula}{}ün sekantın solunda mı yoksa sağında mı
olduğunu gösterir.

Bir tableau dalı hem $\sFmla{\True}{!A}$ hem de
$\sFmla{\False}{!A}$ içeriyorsa \emph{kapalıdır}. Her dal kapalıysa
tableau ağacının tamamına kapalı deriz. Örneğin aşağıdaki, $!A \lor (!B
\land !C) \Sequent (!A \lor !B) \land (!A \lor !C)$ için kapalı bir
tableau örneğidir; yani $\sFmla{\True}{!A \lor (!B \land !C)}$ ve
$\sFmla{\False}{(!A \lor !B) \land (!A \lor !C)}$ ile başlar:
\begin{oltableau}
  [\sFmla{\True}{\formula{A} \lor (\formula{B} \land \formula{C})},
    just=\TAss, checked
    [\sFmla{\False}{(\formula{A} \lor \formula{B}) \land
        (\formula{A} \lor \formula{C})}, just=\TAss, checked
      [\sFmla{\True}{\formula{A}}, just={\TRule{\True}{\lor}[1]}
        [\sFmla{\False}{\formula{A} \lor \formula{B}},
          just={\TRule{\False}{\land}[2]}, checked
          [\sFmla{\False}{\formula{A}},
            just={\TRule{\False}{\lor}[4]}
            [\sFmla{\False}{\formula{B}},
              just={\TRule{\False}{\lor}[4]}, close]
          ]
        ]
        [\sFmla{\False}{\formula{A} \lor \formula{C}},
          just={\TRule{\False}{\land}[2]}, checked
          [\sFmla{\False}{\formula{A}},
            just={\TRule{\False}{\lor}[4]}
            [\sFmla{\False}{\formula{C}},
              just={\TRule{\False}{\lor}[4]}, close]
          ]
        ]
      ]
      [\sFmla{\True}{\formula{B} \land \formula{C}},
        just={\TRule{\True}{\lor}[1]}, checked
        [\sFmla{\False}{\formula{A} \lor \formula{B}},
          just={\TRule{\False}{\land}[2]}, checked
          [\sFmla{\True}{\formula{B}},
            just={\TRule{\True}{\land}[3]}, move by=2
            [\sFmla{\True}{\formula{C}},
              just={\TRule{\True}{\land}[3]}
              [\sFmla{\False}{\formula{A}},
                just={\TRule{\False}{\lor}[4]}
                [\sFmla{\False}{\formula{B}},
                  just={\TRule{\False}{\lor}[4]},close
                ]
              ]
            ]
          ]
        ]
        [\sFmla{\False}{\formula{A} \lor \formula{C}},
          just={\TRule{\False}{\land}[2]}, checked
          [\sFmla{\True}{\formula{B}},
            just={\TRule{\True}{\land}[3]}, move by=2
              [\sFmla{\True}{\formula{C}},
                just={\TRule{\True}{\land}[3]}
                [\sFmla{\False}{\formula{A}},
                  just={\TRule{\False}{\lor}[4]}
                  [\sFmla{\False}{\formula{C}},
                  just={\TRule{\False}{\lor}[4]},close
                ]
              ]
            ]
          ]
        ]
      ]
    ]
  ]
\end{oltableau}

Onay işaretleri, ilgili kuralın bir !!a{formula} için onun altındaki
tüm dallarda uygulanmış olduğunu gösterir. $\otimes$ simgesi bir dalın
kapalı olduğunu gösterir. Kuralların ardından gelen satır numaraları,
kuralın hangi işaretli !!{formula} biçimlerine uygulandığını (yani belirtilen
satırda, aynı dalda bulunan !!{formula} için uygulandığını) gösterir.

Tableau sistemlerinde ispat aramak kolaydır. Bir tableau ağacını aşamalar hâlinde
üretiriz. $0$. aşamada yalnızca başlangıç !!{formula} biçimlerini en üste
yazarız. $n+1$. aşamada her yaprak düğüm için şunları yaparız: Henüz
onay işareti konmamış ilk (en üstteki) işaretli formülü buluruz. Dal
genişletme kuralını bu formüle uygularız. Kural, söz konusu işaretli
formülü içeren her dalda uygulandıktan sonra formüle onay işareti
koyarız. (Yukarıdaki örnek bu algoritma kullanılarak üretilmiştir.)

$\sFmla{\True}{\lforall}$ ve $\sFmla{\False}{\lexists}$ biçimindeki
işaretli !!{formula} biçimlerine hiçbir zaman onay işareti konmaz; dolayısıyla
bunlara özel davranılmalıdır. Bunlar mevcutsa $\TRule{\True}{\forall}$
ve $\TRule{\False}{\lexists}$ kurallarının uygulanabilir tüm terimler
için eninde sonunda uygulanmasını sağlamalıyız. Bu, örneğin her aşamada
bu tür bütün formüllere, her dalda (bu biçimde olmayan en üstteki
işaretli !!{formula} ele alındıktan sonra) kuralları uygulayarak garanti
edilebilir. Kullanılacak $t$ terimi, dalda zaten geçen !!{constant} simgeleri
(hiç !!{constant} geçmiyorsa $\Obj c_0$) ve başlangıçta kullanılan
!!{function} simgeleri---yani başlangıç !!{formula} biçimlerinde
geçenler---kullanılarak kurulabilen ve
karşılık gelen $\sFmla{\True}{!B(t)}$ ya da
$\sFmla{\False}{!B(t)}$ işaretli !!{formula} dalda henüz geçmeyen ilk
$t$ terimidir.

Bu ispat araması kapalı bir tableau üretmezse ya atomik olmayan tüm
!!{formula} biçimlerine onay işareti konmuş sonlu bir açık dal içerir ya da
arama, sonsuz bir dal içermesi gereken sonsuz ve sonlu dallanan bir
ağaç üretir. \olref[cpl]{sec}'de olduğu gibi, dalda
$\sFmla{\True}{!A}$ geçiyorsa $\Sat{M}{!A}$ ve
$\sFmla{\False}{!A}$ geçiyorsa $\Sat/{M}{!A}$ olacak biçimde
!!a{structure}~$\Struct{M}$ tanımlayabiliriz. ($\Theta =
\Setabs{!A}{\sFmla{\True}{!A} \text{ dalda geçer}}$ ve $\Xi =
\Setabs{!A}{\sFmla{\False}{!A} \text{ dalda geçer}}$ alınır.)

Bu ispat arama yöntemiyle üretilen tableau ağaçları, her düğüme bir sekant
atanarak kolayca \Log{G3c} ispatlarına çevrilebilir. Başlangıç
formülleri bir $\Gamma \Sequent \Delta$ sekantına karşılık
geliyorsa başlangıç formüllerinin her birini $\Gamma \Sequent \Delta$
ile etiketleyin. Bir dal, $\sFmla{\True}{!A}$ işaretli !!{formula}
için dal genişletme kuralı uygulanarak genişletilmişse yaprak $!A,
\Pi \Sequent \Lambda$ ile; kural $\sFmla{\False}{!A}$ işaretli
!!{formula} için uygulanmışsa yaprak $\Pi \Sequent \Lambda, !A$ ile
etiketlenir. Tableau'nun bir $\TRule{\True}{\star}$ kuralı uyguladığı
yerde sekanta (ana $!A$ ile) $\LeftR{\star}$ kuralını; bir
$\TRule{\False}{\star}$ kuralı uyguladığı yerdeyse $\RightR{\star}$
kuralını uygulayın. Kuralın öncülleri, tableau ağacına eklenen karşılık gelen
düğümleri etiketleyen sekantlardır. Bir tableau
$\sFmla{\True}{!A}$ ve $\sFmla{\False}{!A}$ ile kapandığında son
yaprak düğümünün etiketi $!A, \Pi \Sequent \Lambda, !A$ biçiminde bir
sekant olacaktır. Tableau ağacında ardışık bazı düğümlere aynı sekant
atanacaktır; yinelenenleri silin. Örneğin yukarıdaki tableau şu
!!{proof} biçimine çevrilir:
\[\small
\Axiom$!A \fCenter !A, !B$
\RightLabel{\RightR{\lor}}
\UnaryInf$!A \fCenter !A \lor !B$
\Axiom$!A \fCenter !A, !C$
\RightLabel{\RightR{\lor}}
\UnaryInf$!A \fCenter !A \lor !C$
\RightLabel{\RightR{\land}}
\BinaryInf$!A \fCenter (!A \lor !B) \land (!A \lor !C)$
\Axiom$!B, !C \fCenter !A, !B$
\RightLabel{\RightR{\lor}}
\UnaryInf$!B, !C \fCenter !A \lor !B$
\RightLabel{\LeftR{\land}}
\UnaryInf$!B \land !C \fCenter !A \lor !B$
\Axiom$!B, !C \fCenter !A, !C$
\RightLabel{\RightR{\lor}}
\UnaryInf$!B, !C \fCenter !A \lor !C$
\RightLabel{\LeftR{\land}}
\UnaryInf$!B \land !C \fCenter !A \lor !C$
\RightLabel{\RightR{\land}}
\BinaryInf$!B \land !C \fCenter (!A \lor !B) \land
(!A \lor !C)$
\RightLabel{\LeftR{\lor}}
\BinaryInf$!A \lor (!B \land !C) \fCenter (!A \lor !B) \land
(!A \lor !C)$
\DisplayProof
\]

\end{document}
